\section{Benchmark on a synthetic Totem market}\label{sec-benchmark}

\subsection{Design}
The synthetic market is built from contributors, so that every mechanism of \cref{sec-anatomy} arises the way it
does in Totem rather than as a shape added to a curve.

\paragraph{Categories.} Four categories with eleven Totem expiries from one month to five years and six dividend
periods (the rest of 2026, then 2027 to 2031); valuation date 14 October 2026.
\begin{center}\input{generated/\lang/tab_categories}\end{center}

\paragraph{Contributors.} The option and dividend panels are drawn independently. Each contributor holds
internally consistent views (\cref{ass-consistent}): dividends with a persistent factor plus per-dividend noise, with a
dispersion that grows with time, and a smooth repo curve (level plus tilt). Option contributors price calls and puts at
five strikes with their own forward, a skewed volatility surface, volatility noise and price noise. Panels shrink with
tenor and are nested.

\paragraph{Consensus and pipeline.} The consensus is the mean after a median/MAD screen; dispersions and counts are
reported as Totem would. The controller pipeline then runs unchanged: parity forwards, allocation on the FO ex-date
profile, bootstrap, covariance, detection.

\paragraph{FO marks.} True dividends with $1\%$ noise and the true repo with a $3$\,bp level error (one standard
deviation).

\paragraph{Scenarios.} Each test name is simulated under every scenario with the same seed, so differences between
scenarios are paired.
\begin{center}\small
\begin{tabular}{p{4.0cm}p{7.4cm}p{3.2cm}}
\toprule
Scenario & Injection & Magnitude\\
\midrule
Spot snapshot mismatch & option contributors price off a different spot & $\pm0.03\%$ to $\pm0.15\%$\\
Dividend bucket error & one period (2027--2029) off in the dividend consensus & $\pm0.15$ to $\pm0.5$ per $100$ of spot\\
Ex-date timing split & $35$--$70\%$ of the option panel moves the dividend closest to an expiry to its other side & one dividend\\
Stale ex-date profile & allocation profile shifted so that one dividend crosses an expiry & $3$--$20$ days past it\\
Forward outlier & consensus calls at one expiry shifted & $3$--$6$ s.d.\ of the carry noise\\
Biased long-dated sub-panel & contributors who remain after the largest panel drop are biased on later dividends & $\pm8\%$ to $\pm20\%$\\
FO repo mismark & FO repo shifted in parallel, Totem untouched & $\pm15$ to $\pm50$\,bp\\
Gross/net bias & every Totem period scaled & $\pm8\%$ to $\pm18\%$\\
\bottomrule
\end{tabular}
\end{center}

\paragraph{Detectors.} (1) Naive gap: largest $|\Delta b|$ over the pillars. (2) Naive shape: largest deviation of the
Totem-implied repo from the interpolation of its two neighbours. (3) B-spline plus Gaussian mixture on the Totem curve:
least-squares cubic B-spline coefficients (six, knots in $\ln T$) clustered by a Gaussian mixture per category with the
number of components chosen by BIC; the score is the negative log-likelihood. (4) The same on the Totem--FO
difference. (5) The structural detector of \cref{sec-detection}.

\paragraph{Protocol.} \vNTrain{} clean names per category calibrate every detector to \vFPR{} false alarms (a quantile
of the score for the baselines, $\kappa$ and the level scale for the structural detector). \vNTestPerCat{} fresh names
per category are then simulated under each of the \vNScen{} scenarios, \vNTest{} test names in total. Because the
design is paired, the signal-to-noise of an injection is computed exactly as $\lambda=\|PL^{-1}(y^{\text{scenario}}-y^{\text{clean}})\|$,
the quantity of \cref{cor-power}.

\subsection{Results}
\begin{table}[ht]
\centering
\caption{Share of names flagged, every detector calibrated to \vFPR{} false alarms on clean history. The clean row is
the realised false-alarm rate on fresh names. For the FO mismark and the gross/net bias a flag is not the desired
outcome (see text). Columns: naive gap, naive shape, B-spline plus Gaussian mixture on the curve and on the Totem--FO
difference, structural detector.}\label{tab-detection}
\input{generated/\lang/tab_detection}
\end{table}

\begin{figure}[ht]
\centering\includegraphics[width=\linewidth]{figures/\lang/detection.pdf}
\caption{Share of names flagged per scenario and detector (the data of \cref{tab-detection}).}
\label{fig-detection}
\end{figure}

\begin{figure}[ht]
\centering\includegraphics[width=\linewidth]{figures/\lang/power.pdf}
\caption{Detection rate against the signal-to-noise ratio $\lambda$ of the injected anomaly, pooled over the six
Totem-side anomalies; $n$ is the number of names per bin. The grey line is the theoretical power of a single known
template at the typical family-wise threshold (\cref{cor-power}).}
\label{fig-power}
\end{figure}

\begin{table}[ht]
\centering
\caption{Area under the ROC curve against the clean names (ranking quality, independent of the threshold).}\label{tab-auc}
\input{generated/\lang/tab_auc}
\end{table}

\begin{table}[ht]
\centering
\caption{Structural detector: share of names flagged per scenario and category.}\label{tab-bycat}
\input{generated/\lang/tab_bycat}
\end{table}

\begin{figure}[ht]
\centering\includegraphics[width=\linewidth]{figures/\lang/attribution.pdf}
\caption{Mechanism named first by the structural detector, per scenario. Level, step, spike and season are the template
families; roughness means the global test fired without a template; no alarm means the name was not flagged.}
\label{fig-attribution}
\end{figure}

\begin{table}[ht]
\centering
\caption{Structural detector: detection, attribution (the first explanation names the injected mechanism and its
location) and FO-mismark flags.}\label{tab-attribution}
\input{generated/\lang/tab_attribution}
\end{table}

\begin{table}[ht]
\centering
\caption{Realised false-alarm rate on fresh clean names, by category (target \vFPR), and the calibrated $\kappa$.}\label{tab-fpr}
\input{generated/\lang/tab_fpr}
\end{table}

\subsection{Reading the results}
\begin{itemize}
\item \textbf{Average detection does not separate the best methods.} Over the six Totem-side anomalies the structural
detector flags \vMeanStruct{} of names, B-spline plus GMM on the difference \vMeanGmmDiff, on the curve \vMeanGmmCurve,
the naive shape check \vMeanShape{} and the naive gap \vMeanNaive.
\item \textbf{Forward outliers} are where carry space and the noise model pay most: \vFwdStruct{} detected against
\vFwdNaive{} for the naive gap, which cannot see a three-to-six standard deviation outlier at a long expiry under the
long-end noise, and \vFwdShape{} for the naive shape check.
\item \textbf{Spot mismatches} (\vSpotStruct) and \textbf{stale ex-date profiles} (\vSeasonStruct) are easy for every
method; the structural detector adds the explanation (\cref{fig-attribution}).
\item \textbf{Dividend bucket errors} (\vBucketStruct{} against \vBucketNaive{} naive) are limited by their size: an
error of $0.15$ to $0.5$ per $100$ of spot is small next to long-dated dividend dispersion, and \cref{fig-power} shows
that the misses are low signal-to-noise cases.
\item \textbf{Ex-date timing splits} are detected in \vTimingStruct{} of names; the naive shape check (\vTimingShape) is
competitive because these spikes are large in repo space.
\item \textbf{A biased long-dated sub-panel} is close to non-identifiable for every method (\vPanelStruct{} structural,
\vPanelNaive{} naive gap): the bias builds up smoothly after the panel drop and mostly projects onto a repo slope. This
is \cref{thm-nonid} at work, not a tuning issue.
\item \textbf{Genuine FO mismarks} are kept apart from bad Totem curves: \vFoStruct{} raise an anomaly (the clean rate)
and \vFoMismark{} are reported as FO-mismark candidates, against \vFoClean{} of clean names. The naive gap and GMM on the
difference raise a bad-curve alarm on \vFoNaive{} and \vFoGmmDiff{} of them. The FO-mismark test finds
\vFoMismarkEuStock{} of the mismarks on EU stocks and \vFoMismarkIndex{} on the index, but only \vFoMismarkHtb{} on
hard-to-borrow names: there contributors' repo views differ by about $60$\,bp, so a $15$--$50$\,bp FO difference is within
the uncertainty of the consensus itself.
\item \textbf{Gross/net bias} is flagged in \vGrossStruct{} of names, mostly as steps at ex-dates for discrete payers;
for an index dividend stream it looks like a repo slope and is reported through $\beta$ instead
(\cref{prop-gross}).
\item \textbf{Power.} Between signal-to-noise $4$ and $6$ the structural detector finds \vPowStructE{} of anomalies,
above $6$ \vPowStructF{} to \vPowStructH, close to the theory of \cref{cor-power}. Below $3$ the naive gap and GMM on the
difference detect more (\vPowGmmDiffC{} against \vPowStructC{} between $2$ and $3$): they also react to the part of an
anomaly that looks like a smooth repo difference, which the structural detector leaves to the FO-mismark test on
purpose. That same property is why they raise alarms on genuine FO mismarks.
\item \textbf{Attribution.} When the structural detector flags a Totem anomaly, its first explanation names the
injected mechanism and location in \vAttr{} of cases.
\end{itemize}

\subsection{Limits of the evidence}
\begin{caveat}
The synthetic market reproduces the mechanisms by construction; it does not show that real Totem data contains only
these. Real data will also carry stale contributors, corporate actions and special dividends, American exercise in
single-stock quotes and data-entry errors. The scale $\kappa$ absorbs a misspecified noise level but not a misspecified
structure, and the magnitudes, while plausible, are not calibrated to any particular history. The benchmark shows that
the method works when the mechanisms are present and is calibrated when they are not; validation on twelve to
twenty-four months of real snapshots with labelled incidents is the next step.
\end{caveat}
